\ifmostrarGabarito
\begin{minted}{python}
from typing import List, Dict, Any

def consolidar_vendas_categoria(
    vendas: List[Dict[str, Any]]
) -> Dict[str, Dict[str, Any]]:
    agrupado = {}
    for pedido in vendas:
        cat = pedido["categoria"]
        val = pedido["valor"]
        if cat not in agrupado:
            agrupado[cat] = {"total": 0.0, "quantidade": 0}
        agrupado[cat]["total"] += val
        agrupado[cat]["quantidade"] += 1
        
    resultado = {}
    for cat, dados in agrupado.items():
        media = dados["total"] / dados["quantidade"]
        resultado[cat] = {
            "total": round(dados["total"], 2),
            "quantidade": dados["quantidade"],
            "media": round(media, 2)
        }
    return resultado
\end{minted}

\textbf{Explicação:}
Realizamos uma passagem sobre os registros acumulando a soma dos valores e a contagem por categoria em um dicionário. Na segunda etapa, calculamos a média dividindo o faturamento pela quantidade de itens.
\else
\linhasResposta{17}
\clearpage
\linhasResposta[FOLHA DE CONTINUACAO / RASCUNHO -- QUESTAO Q11]{30}
\clearpage
\fi
